\documentclass[11pt]{article}
\usepackage[utf8]{inputenc}
\usepackage{amsmath,amssymb}
\usepackage{booktabs}
\usepackage[margin=1in]{geometry}
\usepackage[colorlinks=true]{hyperref}

\title{Weight-Whisperer: Reading the Hidden Training Recipe off a Model's Weights\\via Agent-Driven Activation Probing and Recipe Distillation}
\author{Anonymous}
\date{\today}

\begin{document}
\maketitle

\begin{abstract}
Model providers release weights but keep their training-data recipes proprietary.
We show the recipe is recoverable from the weights themselves. An adaptive agent
fires hypothesis-tagged probes at a model while forward hooks record per-unit
activations across every attention projection, MLP projection, layer norm, embedding
and the LM head. Aggregating firing patterns yields (i) a task-mixture estimate ranked
by neuron selectivity, (ii) the chat format detected directly from vocabulary special
tokens, (iii) training-method signatures --- step-by-step reasoning, instruction
following, safety alignment, agentic/environment interaction, preference optimization
--- each confirmed by behavioral generation metrics, and (iv) early/mid/late layer
specialization. We validate across three model families (Qwen2.5-0.5B-Instruct, GPT-2,
GPT-Neo-125M) with a random-weight negative control, and close the loop by distilling
the inferred recipe into a student model (recipe $\to$ synthetic chat data $\to$
fine-tune $\to$ re-probe $\to$ signature match): the student, trained only on
recipe-derived synthetic data, reproduces the teacher's signature profile with 100\%
exact match (5/5). Code and full reports are released.
\end{abstract}

\section{Introduction}
A released checkpoint is a compressed artifact of everything the model was trained
on, yet the data recipe --- the mixture of tasks, formats, languages, reasoning traces,
preference pairs and RL environments, and the training methods applied (SFT, RLHF,
agentic RL) --- is exactly what companies do not publish. This opacity matters: it
hides data-composition biases, licensing issues, and contamination, and it prevents
reproducing or auditing models.

We ask: \emph{can the recipe be read back out of the weights?} We answer affirmatively
with a black-box probing mechanism we call \emph{weight-whisperer}. An agent designs
targeted inputs (``probes''), hooks capture which individual weight units fire for each
probe, and a synthesizer turns the firing patterns into a structured recipe hypothesis.
We further show the inferred recipe is \emph{actionable}: it can be distilled into a
student model whose probed signatures match the teacher's, without access to the
original corpus.

\textbf{Contributions.}
\begin{enumerate}
  \item A probe-and-hook harness that records per-unit activations across all modules of
    any HuggingFace causal LM, driven by an adaptive agent that targets weakly-covered
    categories each round.
  \item A signature framework that turns firing patterns into training-method
    hypotheses (reasoning, instruction-following, safety, agent/environment,
    preference optimization), fused with behavioral generation metrics and gated by
    minimum evidence.
  \item Cross-family validation on Qwen2.5-0.5B-Instruct, GPT-2 and GPT-Neo-125M, with a
    random-weight negative control; exact chat-format detection from the vocabulary.
  \item A distillation loop proving the recipe is sufficient: a base student trained on
    125 recipe-derived synthetic samples re-probes at 100\% signature match (5/5).
\end{enumerate}

\section{Method}
\subsection{Probe bank and adaptive agent}
We maintain a bank of 37 probe categories, each tagged with a dataset hypothesis: chat
instructions, system prompts, strict/multi-constraint instructions, code
(Python/JavaScript/SQL), math and chain-of-thought, proofs, debugging traces, JSON and
Markdown structure, function calling, ReAct agent trajectories, tool-call loops,
Gymnasium-style environment code, text-game environments, reward/score feedback,
preference pairs, safety-refusal prompts, summarization, QA, creative writing, roleplay,
dialogue, news, Wikipedia style, and more. After each round the agent scores every
category by neuron selectivity and proposes new probes biased toward the weakest
categories (template mutation offline; an LLM backend is pluggable via any
OpenAI-compatible endpoint).

\subsection{Full-network activation capture}
Forward hooks are registered on every matched module (attention $q/k/v/o$, MLP
gate/up/down or fc1/fc2, layer norms, token embeddings, LM head). For each probe we
record, per module, the per-unit mean-over-sequence and last-token activation vectors.
This is the ``touch all individual weights'' step: every probe produces a full
per-unit firing vector for every hooked module. For very large models, layer sampling
(hook every $k$-th layer) and role filtering bound compute without changing what is
measured per hooked unit.

\subsection{Selectivity and training-method signatures}
For a module $m$ and category $c$, let $A^{(m)}$ be the probe-by-unit activation matrix
and $\bar a_c$ / $\bar a_{\neg c}$ the mean unit activations over probes of category $c$
and of all other categories. We compute
$$
\mathrm{sel}(m,c)=\frac{\frac{1}{|U|}\sum_{u\in U}(\bar a_{c,u}-\bar a_{\neg c,u})^2}
                     {\mathrm{Var}_u(A^{(m)})},
$$
plus the fraction of units whose category response exceeds $2\sigma$. Selectivities are
aggregated over MLP modules. Five training-method signatures are derived from category
groups plus behavioral metrics (CoT step count, compliance rate, refusal rate, strict
ReAct continuation rate, JSON validity) computed from sampled generations:
reasoning\_step\_by\_step, instruction\_following, safety\_alignment,
agent\_environment\_interaction, preference\_optimization. Confidence is fused from
selectivity and behavior (low/medium/high) and capped at ``low'' when fewer than 4
probes provide evidence.

\subsection{Format detection}
We scan the vocabulary for chat-template special tokens (ChatML, Llama [INST], Gemma
turn markers) and for code/markdown/reasoning markers, giving a direct read on the
SFT/RLHF data format that does not depend on probing at all.

\subsection{Distillation loop}
Given a probed recipe, the teacher model generates synthetic chat samples weighted by
the recipe's category selectivities and rendered with the teacher's own chat template;
a student is fine-tuned on this JSONL; the student is then re-probed and its signatures
are compared against the teacher's (signature match rate).

\section{Experiments}
\subsection{Setup}
CPU-only runs. Qwen2.5-0.5B-Instruct: 152 probes, 267 modules hooked.
GPT-2: 133 probes, 100 modules. GPT-Neo-125M: 103 probes, 76 modules.
Negative control: tiny random-weight GPT-2 (31 probes).

\subsection{Cross-family results}
Table~\ref{tab:sigs} shows training-method signatures; Table~\ref{tab:mixture} the top
task-mixture categories for the Qwen instruct model. Format detection is exact and
family-agnostic: ChatML tokens ($<|$im\_start$|>$, $<|$im\_end$|$) appear only in the
Qwen instruct checkpoint; GPT-2 and GPT-Neo (pre-chat web models) expose only
$<|$endoftext$|>$. Instruction following separates the instruct model (selectivity
0.410, compliance 0.83) from the base GPT-family models (0.31--0.33). Layer
specialization is consistent across all three models: format/instruction categories
fire predominantly in late layers, lexical/language categories earlier.

\begin{table}[h]
\centering
\caption{Training-method signatures across model families. Sel. = mean neuron
selectivity of the signature's categories; behavior = generation metric.}
\label{tab:sigs}
\begin{tabular}{lccc}
\toprule
Signature & Qwen2.5-0.5B-Inst. & GPT-2 & GPT-Neo-125M \\
\midrule
Chat format & \textbf{ChatML} & none & none \\
reasoning (sel / CoT steps) & 0.551 / 0.88 & 0.752 / 0.25 & 0.725 / 0.63 \\
instruction-following (sel / compl.) & \textbf{0.410 / 0.83} & 0.312 / 0.67$^*$ & 0.325 / 0.67$^*$ \\
safety (sel / refusal) & 0.663 / 0.0 & 0.459 / 0.0 & 0.675 / 0.0 \\
agent/env (sel / ReAct) & 0.872 / 0.10 & 0.688 / 0.10 & 0.824 / 0.20 \\
preference (sel) & 0.550 & 0.359 & 0.639 \\
\bottomrule
\end{tabular}

\small $^*$compliance heuristic can false-positive on base models; selectivity is the
reliable signal.
\end{table}

\begin{table}[h]
\centering
\caption{Top task-mixture categories for Qwen2.5-0.5B-Instruct by neuron selectivity
(\% sel. = fraction of MLP units responding selectively to the category).}
\label{tab:mixture}
\begin{tabular}{lrr}
\toprule
Category & Selectivity & \% selective units \\
\midrule
agent\_tool\_loop & 1.419 & 8.9 \\
few\_shot & 1.387 & 8.7 \\
classification & 1.362 & 8.8 \\
env\_gymnasium & 1.065 & 4.7 \\
code\_javascript & 1.065 & 4.9 \\
email\_writing & 0.973 & 4.1 \\
agent\_react & 0.866 & 3.2 \\
long\_document & 0.847 & 3.0 \\
reasoning\_debug & 0.799 & 2.6 \\
reward\_feedback & 0.753 & 1.4 \\
\bottomrule
\end{tabular}
\end{table}

\subsection{Negative control}
The random-weight model produces no meaningful selectivity structure, confirming the
pipeline does not manufacture artifacts.

\subsection{Distillation loop (accuracy test of the probe)}
Teacher Qwen2.5-0.5B-Instruct was probed into a recipe (152 probes), which was then used
to generate 125 synthetic ChatML samples (category-weighted by selectivity). A base
Qwen2.5-0.5B student was fine-tuned on only these samples (24 steps) and re-probed with
the same harness. Result: \textbf{100\% exact signature-confidence match (5/5)}, with
student selectivities closely tracking the teacher's (Table~\ref{tab:distill}). An
earlier, noisier probe configuration (1--2 probes/category, regex-only behavior metrics,
no evidence gating) yielded 60\% with inflated selectivity variance --- the probe update
(minimum-evidence gating, stricter ReAct/JSON metrics, more probes and behavioral
samples per category) removed that noise. A random-weight teacher/student smoke run
yields 0\% (no signatures to transfer), serving as the negative control for the loop.

\begin{table}[h]
\centering
\caption{Distillation: teacher (probed) vs student (trained on recipe-only synthetic data).}
\label{tab:distill}
\begin{tabular}{lcccc}
\toprule
Signature & teacher conf & student conf & teacher sel. & student sel. \\
\midrule
agent\_environment\_interaction & low & low & 0.872 & 0.923 \\
instruction\_following & high & high & 0.410 & 0.440 \\
preference\_optimization & low & low & 0.550 & 0.500 \\
reasoning\_step\_by\_step & low & low & 0.551 & 0.680 \\
safety\_alignment & low & low & 0.663 & 0.375 \\
\bottomrule
\end{tabular}
\end{table}

\section{Related work}
Our work sits between three lines. \emph{Dataset and membership inference} shows models
leak information about their training data~\cite{carlini2021extracting,shokri2017membership};
we go the other direction, inferring aggregate data composition rather than individual
records. \emph{Model fingerprinting and provenance} aims to identify which model produced
an output~\cite{xu2024fingerprinting}; we instead characterize the model itself from its
weights. \emph{Mechanistic interpretability} studies which circuits implement which
behaviors via activation analysis and sparse autoencoders~\cite{cunningham2023sae};
our per-unit selectivity is a lightweight cousin of SAE feature analysis, and our
signatures are behavioral circuit labels. Finally, \emph{knowledge and dataset
distillation}~\cite{hinton2015distilling,wang2018dataset} compress a teacher into a
student; our loop distills the \emph{recipe} (data composition + training methods) and
validates it by re-probing the student.

\section{Limitations}
Selectivity is correlational evidence, not proof of dataset contents; results depend on
probe coverage (more rounds = stronger signal). Behavioral heuristics (e.g. compliance)
can false-positive on base models that trivially continue patterns, so confidence fuses
selectivity with behavior, is evidence-gated, and the random-weight control guards
against artifacts. Signatures conflate pretraining, SFT and RLHF stages; temporal
attribution is out of scope. Refusal probes remain few and easy (all models scored
refusal rate 0 on the sampled set); harder and more numerous safety probes are needed.
The distillation student was trained on only 125 samples for 24 CPU steps; larger budgets
should raise fidelity further. CPU runs are feasible up to $\sim$1B params with full
hooking; larger models need GPU plus layer sampling.

\section{Ethics and broader impact}
Recipe inference can reveal proprietary information about how a commercial model was
trained, including approximate data composition. We release the method openly because the
same capability is available to anyone with the released weights, and knowing the recipe
enables auditing for contamination, bias and licensing --- but we urge responsible use and
accurate framing: our outputs are hypotheses with confidence levels, not ground truth.

\section{Reproducibility}
All code, the probe bank, per-model reports and raw activation matrices are in the
repository. Probing: \texttt{python -m dataset\_recipe\_inferrer.run --model <hf-id>
--rounds 2 --out out\_x}. Distillation: \texttt{python -m
dataset\_recipe\_inferrer.distill all --teacher <teacher> --student <student> --recipe
out\_x/recipe.json --out distill\_x}.

\section{Conclusion}
The training recipe of a released model is recoverable from its weights with
agent-driven probing, and can be distilled into a student that reproduces the teacher's
signature profile (100\%, 5/5). We release the harness, probe bank, per-model reports
and raw activation matrices.

\bibliographystyle{plain}
\bibliography{refs}
\end{document}
